\section{Luồng vận hành}\label{sec:4}


Một tác vụ chạy các pha theo thứ tự CHECK\_CFG → LOAD\_BIAS (nếu bật hậu xử lý) → LOAD\_A → LOAD\_B → COMPUTE → STORE\_C. Mỗi lúc chỉ một pha chạy, và mỗi pha có đúng một khối làm chính. Pha sau chỉ bắt đầu khi pha trước thành công.


\subsection{Các pha của một tác vụ}\label{sec:4.1}

\begin{spectable}{|L{0.14}|L{0.13}|L{0.53}|L{0.2}|}{Các pha}{tab:4-1}{\textbf{Pha} & \textbf{Khối chính} & \textbf{Làm gì} & \textbf{Xong khi}}
CHECK\_CFG & Main FSM & Kiểm tra kích thước, căn địa chỉ và các trường dự trữ theo điều kiện hợp lệ của thanh ghi. Sai: lỗi mã 1; không truy cập bus hoặc scratchpad. & Kiểm tra xong \\\hline

LOAD\_BIAS & DMA & Chép N word bias từ DRAM vào buffer BIAS. Chỉ có khi \sig{POST_EN} = 1 & SRAM nhận word cuối \\\hline

LOAD\_A & DMA & Chép A (M hàng × K/4 word) vào buffer A & SRAM nhận word cuối \\\hline

LOAD\_B & DMA & Chép B (K hàng × N/4 word) vào buffer B & SRAM nhận word cuối \\\hline

COMPUTE & Array, VPU & Array tính mọi tile; mỗi hàng kết quả đi qua VPU (hậu xử lý hoặc đi thẳng) rồi ghi vào buffer C & \sig{pe_done} \\\hline

STORE\_C & DMA & Chép C (M·N word) hoặc Y (M·N/4 word) từ buffer C ra DRAM & Phản hồi B cuối là OKAY \\\hline

\end{spectable}



\subsection{Luồng một tác vụ}\label{sec:4.2}



Có ba mức “xong”: \sig{o_done} của DMA báo xong một lệnh DMA; \sig{pe_done} báo toàn bộ kết quả đã nằm trong buffer C; STATUS.DONE báo toàn bộ kết quả đã nằm trong DRAM. Chỉ mức cuối cùng hiện ra cho CPU.

\subsection{Quyền truy cập scratchpad theo pha}\label{sec:4.3}

\begin{spectable}{|L{0.18}|L{0.2}|L{0.2}|L{0.2}|L{0.22}|}{Ai truy cập buffer trong mỗi pha}{tab:4-2}{\textbf{Pha} & \textbf{A} & \textbf{B} & \textbf{BIAS} & \textbf{C}}
IDLE, CHECK\_CFG & — & — & — & — \\\hline

LOAD\_BIAS & — & — & DMA ghi & — \\\hline

LOAD\_A & DMA ghi & — & — & — \\\hline

LOAD\_B & — & DMA ghi & — & — \\\hline

COMPUTE & Array đọc & Array đọc & VPU đọc & VPU ghi \\\hline

STORE\_C & — & — & — & DMA đọc \\\hline

\end{spectable}


\paragraph{Quyền truy cập}\label{tab:4-3}
\begin{itemize}
\item Mỗi buffer chỉ có tối đa một bên truy cập trong mỗi pha, theo các quy tắc địa chỉ và giao tiếp bên dưới. Main FSM bảo đảm điều này bằng cách chạy tuần tự các pha. Vì vậy mỗi bank vật lý chỉ cần SRAM một cổng
\item Module SRAM có assertion (mô phỏng) phát hiện truy cập sai bên trong một pha
\end{itemize}


\subsection{Xử lý lỗi}\label{sec:4.4}


\paragraph{Lỗi}\label{tab:4-4}
\begin{itemize}
\item Cấu hình sai: CHECK\_CFG kết thúc tác vụ với \sig{ERR_CODE} = 1; không có giao dịch AXI nào, không truy cập scratchpad
\item Lỗi AXI trong một lệnh DMA: DMA ngừng phát burst mới, nhận đủ beat R hoặc phản hồi B của burst đã phát, rồi báo \sig{o_error}. Main FSM kết thúc tác vụ với \sig{ERR_CODE} = 2 (LOAD\_*) hoặc 3 (STORE\_C) và không phát lệnh nào nữa
\item Array và VPU không có đầu ra lỗi
\end{itemize}


\subsection{Reset và soft reset}\label{sec:4.5}


\paragraph{Reset}\label{tab:4-5}
\begin{itemize}
\item \sig{rst_n} tích cực mức thấp, bất đồng bộ. \sig{gemm_accel_top} có bộ đồng bộ 2 flip-flop (bật bất đồng bộ, nhả đồng bộ với \sig{clk}); đầu ra \sig{rst_n_sync} cấp cho mọi module
\item \sig{rst_n_sync} = 0 đưa mọi trạng thái về giá trị reset. Nội dung scratchpad không bị xóa, coi là không xác định
\item Soft reset không đi vào chân reset; nó là xung xóa đồng bộ \sig{soft_clr} một chu kỳ do Main FSM phát tới Register File, DMA, SRAM, Array, VPU
\item Phạm vi soft reset: FSM, DMA, Array, VPU, phần điều khiển SRAM; STATUS và \sig{PERF_*}. Thanh ghi R/W giữ nguyên. Nội dung scratchpad không xác định
\item Nếu DMA đang chạy lệnh: FSM bật \sig{i_dma_abort}; DMA không phát burst mới, hoàn tất mọi burst đã cam kết (đã bật VALID hoặc đã gửi beat W), rồi kết thúc lệnh bằng \sig{o_done} hoặc \sig{o_error}. Sau đó FSM phát \sig{soft_clr}
\item Nếu DMA không chạy lệnh (IDLE, CHECK\_CFG, COMPUTE): FSM phát \sig{soft_clr} ở chu kỳ kế tiếp; COMPUTE đang chạy bị hủy
\item Trong cả chuỗi soft reset: BUSY = 1; START và SOFT\_RESET ghi lặp lại bị bỏ qua
\item Sau \sig{soft_clr}: STATUS = 0; tác vụ bị hủy không báo DONE cũng không báo ERROR; vùng đích trong DRAM không xác định
\item Nếu yêu cầu soft reset đến đúng chu kỳ DMA kết thúc lệnh, coi như DMA đã rảnh và áp dụng (Nếu DMA không chạy lệnh (IDLE, CHECK\_CFG, COMPUTE): FSM phát \sig{soft_clr} ở chu kỳ kế tiếp; COMPUTE đang chạy bị hủy ). Khi đã có yêu cầu soft reset, FSM không phát lệnh mới
\end{itemize}


% Mục 7 — Thành phần bổ sung (nếu còn thời gian)
